\documentclass[11pt,reqno]{amsart}

\usepackage[margin=1.15in]{geometry}
\usepackage{amsmath,amssymb,amsthm}
\usepackage{booktabs}
\usepackage{hyperref}

\hypersetup{
    colorlinks=true,
    linkcolor=blue!60!black,
    citecolor=blue!60!black,
    urlcolor=blue!60!black
}

\newtheorem{theorem}{Theorem}[section]
\newtheorem{lemma}[theorem]{Lemma}
\newtheorem{proposition}[theorem]{Proposition}
\newtheorem{corollary}[theorem]{Corollary}
\theoremstyle{remark}
\newtheorem{remark}[theorem]{Remark}

\newcommand{\ES}{\mathcal{E}_S}
\newcommand{\TV}{\mathrm{TV}}
\newcommand{\op}{\mathrm{op}}
\newcommand{\SRB}{\mathrm{SRB}}
\newcommand{\loc}{\mathrm{loc}}
\newcommand{\two}{\mathrm{two}}
\newcommand{\defect}{\mathrm{defect}}
\newcommand{\diam}{\mathrm{diam}}
\newcommand{\supp}{\mathrm{supp}}
\newcommand{\Var}{\mathrm{Var}}
\newcommand{\Cov}{\mathrm{Cov}}

\title{Non-Asymptotic Symbolic Ergodeviation}
\author{Fyyre (James) \\ \texttt{fyyre@fyyre.net}}
\date{\today}

\begin{document}

\begin{abstract}
Classical dynamical and information-theoretic invariants---including Shannon entropy rate, Kolmogorov--Sinai entropy, topological entropy, Lyapunov exponents, correlation dimension, and Lempel--Ziv complexity---either require asymptotic infinite-horizon limits or fail to distinguish benign local metric jitter from structured non-ergodic drift on finite windows. We specify and analyze a finite-window composite functional, \emph{Symbolic Ergodeviation} ($\ES$), combining the $p$-Wasserstein distance between empirical and reference $\ell$-word marginals with a one-sided Dirichlet energy defect of the reference Rauzy graph's Fiedler vector. We prove that natural two-sided Fiedler ratios and empirical visit-graph Laplacians fail on finite orbits, establish exact zero-set, local-jitter stability, and non-asymptotic Azuma concentration bounds under Dobrushin contraction, bound the spectral perturbation of estimated baselines via the Davis--Kahan $\sin\theta$ theorem, and certify exact rational primal--dual separation on a finite-state Markov chain.
\end{abstract}

\maketitle

\section*{Status}

Symbolic ergodeviation, written $\ES$, is not a classical named invariant. Shannon entropy rate, Kolmogorov--Sinai entropy, topological entropy, Lyapunov exponents, correlation dimension, and Lempel--Ziv complexity are. Each of those is the wrong functional for the distinction drawn below, and Section~\ref{sec:ties} records explicit pairs on which they tie while a transport cost does not.

The ingredients used here are standard: a physical (SRB) measure and its symbolic marginal, the $1$-Wasserstein metric, and the algebraic connectivity of a weighted Rauzy graph. What is specified and proved is one finite-window composite of those ingredients. The proofs below apply to that composite. They do not create a uniqueness theorem for the formula, and they do not assert that an operating-system kernel or a packet trace is an Axiom~A attractor.

A second formula, the two-sided ratio of Fiedler energies, looks equally plausible from the phrase ``penalized by algebraic connectivity'' and is rejected in Section~\ref{sec:bad-penalties}. On a deterministic rare orbit it shrinks a transport cost of $4/5$ down to $4/17$.

Population values in the worked example are exact rational numbers, certified by an explicit coupling and a $1$-Lipschitz dual witness. Monte Carlo numbers are labeled as such and come from \texttt{scratch/verify\_es.py}.

\section{Claim}\label{sec:claim}

Let a baseline process be a finite-state exponentially mixing shift, with law $\nu$, thought of as the symbolic image of a physical measure when a smooth model exists. Let $\hat\mu_{N,\ell}$ be the empirical law of length-$\ell$ words in a window of $N$ symbols. Let $\lambda_2$ be the algebraic connectivity of the flux-weighted graph of transitions that $\nu$ allows, and let $\hat\Delta_N$ be the Dirichlet energy of a matching Fiedler vector along the window.

For a fixed weight $\eta \ge 0$ and a block length $\ell \ge 2$,
\begin{equation}\label{eq:es-def}
\ES
=
W_p(\hat\mu_{N,\ell}, \nu_\ell)
\left(
1 + \eta \frac{(\lambda_2 - \hat\Delta_N)_+}{\lambda_2}
\right).
\end{equation}

The theorems are:
\begin{enumerate}
\item $\ES = 0$ if and only if the window's $\ell$-marginal equals $\nu_\ell$. The connectivity factor does not move the zero set.
\item Under the baseline, $\ES \to 0$ almost surely, and a fully explicit non-asymptotic tail holds for every finite $N$ (Azuma, with a variance bound on the expectation). The constant is valid and not sharp.
\item Moving every symbol by at most $\eta_{\loc}$ in the symbol metric raises $\ES$ by at most the factor $(1+\eta)$ times a term linear in $\eta_{\loc}$. High entropy inside a small metric ball stays cheap.
\item Any alternative whose $\ell$-marginal is at least $\tau$ away in $W_p$ scores at least $\tau$, whether or not its entropy rate has moved. If it also realizes only a fraction of the baseline Fiedler energy, the score is strictly larger than the pure transport cost.
\item Shannon rate, Kolmogorov--Sinai entropy, Lempel--Ziv rate, and correlation dimension can be held fixed while $W_1$, hence $\ES$, is arbitrarily large. The Lyapunov exponent of the doubling map is the same on Lebesgue and on the fixed point.
\end{enumerate}

\section{Assumptions}\label{sec:assumptions}

\noindent\textbf{Alphabet and metric.} $A$ is finite, $|A| = m$. A metric $\delta$ on $A$ is fixed. The discrete metric $\delta(a,b) = 1_{a \ne b}$ is allowed; the stability statement that local jitter is cheap needs $\delta$ to come from the observation geometry (bins that sit next to each other in the counter space are close). On words of length $\ell$,
\[
d_\theta(u,v) = \sum_{j=0}^{\ell-1} \theta^j \, \delta(u_j, v_j), \qquad \theta \in (0,1).
\]
Write $D = \diam(A^\ell, d_\theta) < \infty$. If $\delta$ is a metric then so is $d_\theta$.

\medskip\noindent\textbf{Baseline.} $(s_n)_{n \ge 0}$ is a stationary ergodic Markov chain on $A$ with transition matrix $P$ and stationary row $q$. The chain is irreducible and aperiodic on its support. The Dobrushin coefficient
\[
\kappa(P) = \max_{a,b \in A} \|P(a,\cdot) - P(b,\cdot)\|_{\TV}
\]
satisfies $\kappa < 1$, where $\|\mu-\nu\|_{\TV} = \tfrac12 \|\mu-\nu\|_1$. Then $\|\mu P - \nu P\|_{\TV} \le \kappa \|\mu-\nu\|_{\TV}$. If only some power contracts, $\kappa(P^m) = \kappa_m < 1$, the same argument applies with the constant recorded in Remark~\ref{rem:m-step}.

The reference marginal $\nu_\ell$ is the law of $(s_0,\ldots,s_{\ell-1})$. The transition tensor of an order-$1$ chain is the case $\ell = 2$: $\nu_2(a,b) = q_a P_{ab}$.

\medskip\noindent\textbf{Where SRB enters, and where it does not.} If the baseline is a $C^{1+\alpha}$ Axiom~A attractor and the partition is a Markov partition, Bowen's theorem supplies a unique physical measure $\mu_{\SRB}$, invariant, with conditional measures absolutely continuous on unstable manifolds, and the symbolic process is a H\"older Gibbs state for $\varphi = -\log |\det Df|_{E^u}|$. Exponential mixing of H\"older observables gives some $m$ with $\kappa_m < 1$. In that case $\nu = \pi_* \mu_{\SRB}$. The measure of maximal entropy is the same object only when $\varphi$ is cohomologous to a constant. Nothing below requires that identification.

A performance-counter trace or a packet trace does not inherit an SRB theorem from this citation. For those series, $\nu$ is the law of a trusted baseline corpus. Every theorem in Sections~\ref{sec:elementary} and~\ref{sec:finite-sample} is a theorem about that $\nu$. The name ``SRB'' is reserved for the smooth case just stated.

\section{What the Classical Invariants Tie, and What They Do Not}\label{sec:ties}

\subsection{An entropy-blind pair}\label{subsec:entropy-blind}

Let $M \ge 2$ and let $A = \{0,1,\ldots,M\}$ with $\delta(i,j) = |i-j|$. Let $\rho_0$ be i.i.d.\ uniform on $\{0,1\}$ and $\rho_1$ i.i.d.\ uniform on $\{M-1, M\}$.

Both processes have Shannon entropy rate $1$~bit. Both are Bernoulli, so the Kolmogorov--Sinai entropy of the shift under each measure is $1$~bit. The Lempel--Ziv compression rate of a stationary ergodic source converges to its entropy rate, so both have Lempel--Ziv rate $1$~bit. Each support is two points, so the correlation dimension is $0$. The topological entropy of each support subshift is $\log 2$, read off the language only in the limit of block length.

The identity coupling that sends $0 \mapsto M-1$ and $1 \mapsto M$ has cost $M-1$. The function $f(i) = i$ is $1$-Lipschitz and
\[
\int f \,d\rho_1 - \int f \,d\rho_0 = \tfrac12\bigl((M-1)+M\bigr) - \tfrac12\bigl(0+1\bigr) = M-1.
\]
Kantorovich--Rubinstein gives $W_1(\rho_0,\rho_1) = M-1$, which grows without bound as the two pairs are moved apart, with every entropy-type number held fixed.

On a finite metric space $W_p \ge W_1$. Whenever $\ES$ is defined for a baseline equal to one of these laws, the factor in parentheses is at least $1$, so the same pair is separated by $\ES$ by at least $M-1$.

\subsection{A Lyapunov-blind pair}\label{subsec:lyapunov-blind}

On the doubling map $T(x) = 2x \bmod 1$, one has $|T'| = 2$ at every point, so every orbit has Lyapunov exponent $\log 2$. Lebesgue measure is the SRB measure. Its symbolic image under the partition $\{[0,1/2),[1/2,1)\}$ is fair Bernoulli, whose Kolmogorov--Sinai entropy equals $\log 2$ by Pesin's identity. The Dirac mass at the fixed point $0$ is invariant, has entropy $0$, and is coded by the constant sequence $000\ldots$. With the discrete symbol metric and $\ell = 1$,
\[
W_1(\delta_0, \tfrac12\delta_0 + \tfrac12\delta_1) = \tfrac12.
\]
The exponent that numerical Lyapunov estimators target is the same on the two measures. Entropy sees this particular pair. Section~\ref{subsec:entropy-blind} is the pair entropy does not see. The point of listing both is that no one of these invariants is a finite-window transport distance to a reference marginal.

\subsection{Asymptotic grammar}

Topological entropy is $\lim_{n\to\infty} n^{-1} \log \#\{\text{words of length }n\}$. It is a property of the support language along an infinite orbit. A window of length $N$ does not determine it, and two measures with the same support share it. It does not define a distance from $\hat\mu_{N,\ell}$ to $\nu_\ell$.

\section{The Graph Whose Connectivity Is Used}\label{sec:graph}

Let $\ell \ge 2$ and $k = \ell - 1$. Vertices are the $k$-words in the support of $\nu_k$. For an order-$1$ chain, $k = 1$: vertices are symbols.

The forbidden set is the set of $(k+1)$-words with $\nu$-mass zero. The Rauzy graph keeps the complementary edges. Give the undirected edge $\{u,v\}$ the stationary flux
\[
w(u,v) = \xi(u \to v) + \xi(v \to u),
\]
where $\xi(a,b) = q_a P_{ab}$ in the order-$1$ case, and the analogous cylinder mass in general. Self-loops do not appear in $w$.

The weighted Laplacian is $(L\psi)_u = \sum_v w(u,v)(\psi_u - \psi_v)$. Its kernel is the constants precisely because the chain is irreducible, so the underlying graph is connected. Algebraic connectivity $\lambda_2 > 0$ is the second-smallest eigenvalue (Fiedler). The Rayleigh theorem gives, for a unit Fiedler vector $\psi \perp \mathbf{1}$,
\[
\lambda_2 = \psi^\top L \psi = \sum_{\{u,v\}} w(u,v)\,(\psi_u - \psi_v)^2.
\]
The same quadratic form, written on ordered transitions, is the baseline energy
\[
\Delta_* = \sum_{u,a} \xi(u,a)\,(\psi(u)-\psi(ua))^2.
\]
Each undirected edge contributes once through $w$, and $\xi(\cdot\to\cdot)+\xi(\cdot\leftarrow\cdot) = w$, so
\[
\Delta_* = \lambda_2.
\]
If the eigenvalue has multiplicity $m_2 > 1$, replace $(\psi_u-\psi_v)^2$ by $\|\Psi(u)-\Psi(v)\|^2$ for an orthonormal basis of the eigenspace. The identity becomes $\Delta_* = m_2 \lambda_2$, which is basis-independent, and the ratio $(\Delta_* - \hat\Delta)_+/\Delta_*$ is unchanged in form. The worked example has multiplicity one.

\medskip\noindent\textbf{Empirical energy.} On a window $s_0,\ldots,s_{N-1}$, with $N' = N-\ell+1$ consecutive $\ell$-words,
\[
\hat\Delta_N = \frac{1}{N'}\sum_{t=0}^{N'-1} c(s_t,\ldots,s_{t+\ell-1}),
\]
where $c = (\psi(u)-\psi(u'))^2$ if the transition is legal, and $c = 0$ if the word is forbidden. Legal mass that fails to cross a Fiedler gap, and illegal mass, both lower $\hat\Delta_N$. Illegal mass is also charged by $W_p$, because $\nu$ gives it nothing to sit on.

Under a stationary draw from $\nu$, every term has mean $\Delta_*$, so $\mathbb{E}[\hat\Delta_N] = \lambda_2$.

\section{Definition}\label{sec:def}

Let $\hat\mu_{N,\ell}$ be the empirical measure on $\ell$-words in the window. For $p \ge 1$ and $\eta \ge 0$,
\[
\defect(\hat\Delta_N)
=
\frac{(\lambda_2 - \hat\Delta_N)_+}{\lambda_2},
\qquad
\ES
=
W_p(\hat\mu_{N,\ell}, \nu_\ell)
\bigl(1 + \eta \,\defect(\hat\Delta_N)\bigr).
\]
The defect lies in $[0,1]$, so the factor lies in $[1, 1+\eta]$. In particular
\[
W_p(\hat\mu_{N,\ell}, \nu_\ell)
\le
\ES
\le
(1+\eta)\, W_p(\hat\mu_{N,\ell}, \nu_\ell).
\]
Setting $\eta = 0$ returns plain Wasserstein. The value $\eta = 1$ used in Section~\ref{sec:worked} weights a total collapse of Fiedler energy equally with the transport cost: the factor equals $2$ when $\hat\Delta_N = 0$, and equals $1$ whenever the window meets or exceeds the baseline energy.

On a space of diameter $D$, $p \le q$ implies $W_p \le W_q \le D^{1-p/q} W_p^{p/q}$. Bounds proved for $W_1$ transfer to $W_p$ up to factors that depend only on $D$ and $p$. The concentration proof is written for $p = 1$.

\section{Two Penalties That Do Not Work}\label{sec:bad-penalties}

\subsection{Algebraic connectivity of the visit graph}

Let $H_N$ be the unweighted graph whose vertices are the $k$-words actually seen, with an edge for each consecutive pair actually seen. A window that has not repeated a state induces a path $P_n$ on $n \le N$ vertices.

\begin{lemma}[Path test vector]
$\lambda_2(P_n) \le 12/(n(n+1))$.
\end{lemma}
\begin{proof}
Let $v_j = j - (n+1)/2$ for $j = 1,\ldots,n$. Then $\sum v_j = 0$ and
\[
\|v\|_2^2 = \frac{n(n^2-1)}{12}.
\]
Adjacent coordinates differ by $1$, and there are $n-1$ edges, so $v^\top L v = n-1$. Rayleigh's theorem yields
\[
\lambda_2 \le \frac{n-1}{n(n^2-1)/12} = \frac{12}{n(n+1)}.\qedhere
\]
\end{proof}

Thus $\lambda_2(H_N) \to 0$ along any non-repeating window, including a window drawn from the baseline on a large alphabet, long before the cover time. A factor $1/\lambda_2(H_N)$ diverges on typical finite orbits. That factor is not a non-asymptotic penalty.

The Dirichlet form $\hat\Delta_N$ avoids this. It is an average of a fixed bounded function of the chain, and it is unbiased for $\lambda_2$ at every $N$, before cover.

\subsection{The two-sided ratio}\label{subsec:two-sided}

The ratio
\[
\Pi^{\two} = \frac{\lambda_2(1+\varepsilon)}{\hat\Delta_N + \varepsilon\lambda_2}, \qquad \varepsilon \in (0,1],
\]
is near $1$ when $\hat\Delta_N$ is near $\lambda_2$, and it is large when $\hat\Delta_N$ is near $0$. It is also strictly smaller than $1$ whenever $\hat\Delta_N > \lambda_2$. Any alternative that crosses the Fiedler cut more than the baseline does is discounted.

Section~\ref{sec:worked} computes a deterministic alternation on the rare cut of a $3$-state chain: $W_1 = 4/5$, $\hat\Delta = 2$, $\lambda_2 = 1/2$. With $\varepsilon = 1/4$,
\[
\Pi^{\two} \cdot \tfrac45 = \tfrac{4}{17} \approx 0.2353.
\]
The one-sided definition scores the same orbit at $4/5$. In an $800$-window Monte Carlo of the baseline at $N = 400$ (Section~\ref{sec:worked}), the two-sided baseline scores had maximum $0.202$. The anomaly is still above that maximum, and the gap has collapsed from several tenths to about $0.03$. The one-sided score $4/5$ remains well above the one-sided baseline maximum $0.236$ from the matching experiment.

\begin{lemma}[No-hiding]
If the prefactor is at least $1$ for every window, then every alternative with $W_p \ge \tau$ has $\ES \ge \tau$. The two-sided ratio violates the hypothesis. The one-sided factor $1+\eta\,\defect$ satisfies it.
\end{lemma}

\section{Elementary Properties}\label{sec:elementary}

\begin{proposition}[Zero set]\label{prop:zero}
For a population measure $\rho$ on $\ell$-words, write $\Delta(\rho)$ for the legal Dirichlet form and
\[
\mathcal{E}(\rho) = W_p(\rho, \nu_\ell)\bigl(1+\eta\,\defect(\Delta(\rho))\bigr).
\]
Then $\mathcal{E}(\rho) = 0$ if and only if $\rho = \nu_\ell$.
\end{proposition}
\begin{proof}
The prefactor is at most $1+\eta$ and at least $1$. The product vanishes if and only if $W_p = 0$. Since $d_\theta$ is a metric, $W_p = 0$ if and only if the two probabilities agree. If $\rho = \nu_\ell$ then $\Delta(\rho) = \Delta_* = \lambda_2$, so the defect is zero, which matches.
\end{proof}

\begin{proposition}[Almost-sure limit]\label{prop:as}
For $\nu$-almost every baseline orbit, $\ES \to 0$.
\end{proposition}
\begin{proof}
The shift is ergodic and each cylinder indicator is integrable. Birkhoff's theorem gives $\hat\mu_{N,\ell} \to \nu_\ell$ almost surely in total variation, hence in $W_p$, on a finite set. The same theorem applied to the bounded function $c$ gives $\hat\Delta_N \to \lambda_2$, so the defect tends to $0$ and the prefactor tends to $1$.
\end{proof}

\begin{proposition}[Jitter]\label{prop:jitter}
Let $K$ be a kernel on $A$ with $W_1(\delta_a, K(a,\cdot)) \le \eta_{\loc}$ for every $a$, the Wasserstein cost being computed with $\delta$. Let $\rho$ be any law on $A^{\mathbb{N}}$, and let $\rho^K$ be the law after applying $K$ independently at each coordinate. Then
\[
W_1(\rho^K_\ell, \nu_\ell)
\le
W_1(\rho_\ell, \nu_\ell) + \eta_{\loc}\frac{1-\theta^\ell}{1-\theta},
\]
and therefore
\[
\mathcal{E}(\rho^K_\ell)
\le
(1+\eta)\left(
W_1(\rho_\ell, \nu_\ell) + \eta_{\loc}\frac{1-\theta^\ell}{1-\theta}
\right).
\]
\end{proposition}
\begin{proof}
Couple each clean symbol to its noisy image at cost $\le \eta_{\loc}$, independently. The resulting word coupling has $d_\theta$-cost at most $\eta_{\loc} \sum_{j=0}^{\ell-1} \theta^j$. That is the first inequality, after the triangle inequality through $\rho_\ell$. The second is the comparison $\mathcal{E} \le (1+\eta) W_1$, which holds pathwise because the defect is at most $1$. The prefactor cannot amplify local noise by more than $1+\eta$.
\end{proof}

Entropy of $\rho^K$ can be large. The bound does not depend on that entropy. What it depends on is $\eta_{\loc}$ in the symbol metric. With the discrete metric, a flip to any other bin costs $1$, and the proposition correctly refuses to call that flip local.

\begin{proposition}[Structured drift]\label{prop:drift}
Let $\rho$ be a probability on $\ell$-words with $W_p(\rho, \nu_\ell) \ge \tau$. Then $\mathcal{E}(\rho) \ge \tau$. If also $\Delta(\rho) \le \alpha \lambda_2$ for some $\alpha \in [0,1)$, then
\[
\mathcal{E}(\rho) \ge \tau \bigl(1 + \eta(1-\alpha)\bigr).
\]
\end{proposition}
\begin{proof}
The prefactor is at least $1$, and it is at least $1+\eta(1-\alpha)$ when the energy shortfall is at least $1-\alpha$.
\end{proof}

For a periodic sequence of period $p$, the empirical $\ell$-marginal and $\hat\Delta_N$ differ from the periodic occupation by at most $O((p+\ell)/N)$ in total variation and in the bounded energy. For all large $N$, $\ES$ exceeds half the population value of that periodic occupation.

\section{Finite-Sample Bound}\label{sec:finite-sample}

Throughout this section the window is stationary under $\nu$, $\kappa = \kappa(P) < 1$, $p = 1$, and $N' = N - \ell + 1 \ge 1$.

\begin{lemma}[One-symbol Lipschitz constant of $W_1$]\label{lem:lip}
Changing a single symbol $s_j$ changes $W_1(\hat\mu_{N,\ell}, \nu_\ell)$ by at most
\[
c = \frac{D \,\ell}{N'}.
\]
\end{lemma}
\begin{proof}
At most $\ell$ windows cover site $j$. Each of those windows contributes a count of $1$. Replacing it moves one count from one word to another, so the $\ell_1$ change in the count vector is at most $2\ell$. Hence $\|\hat\mu-\hat\mu'\|_1 \le 2\ell/N'$ and the total-variation distance is at most $\ell/N'$. On a space of diameter $D$, $W_1 \le D \,\|\cdot\|_{\TV}$, because a maximal coupling realizes $\mathbb{P}(X \ne Y) = \|\cdot\|_{\TV}$ and $d_\theta \le D$.
\end{proof}

\begin{lemma}[Influence]\label{lem:influence}
Let $f$ be a function of $(s_0,\ldots,s_{N-1})$ with $|f(s)-f(s')| \le \sum_{i: s_i \ne s'_i} c_i$. The Doob increments of $f$ along the chain satisfy
\[
\bigl|\mathbb{E}[f \mid s_0,\ldots,s_j] - \mathbb{E}[f \mid s_0,\ldots,s_{j-1}]\bigr|
\le
\sum_{t \ge 0} c_{j+t}\, \kappa^{t}.
\]
\end{lemma}
\begin{proof}
Fix the past up to time $j-1$. For two possible values $a,b$ of $s_j$, couple the two futures by the maximal coupling so that $\mathbb{P}(s_{j+t}^{(a)} \ne s_{j+t}^{(b)}) \le \kappa^t$. This is the Dobrushin contraction iterated. The conditional expectations $g(a), g(b)$ therefore differ by at most $\sum_{t \ge 0} c_{j+t} \kappa^t$. The random variable $g(s_j)$ and its conditional mean both lie between $\min g$ and $\max g$, so the Doob increment is at most that oscillation.
\end{proof}

\begin{lemma}[Variance of a word frequency]\label{lem:var}
Let $J_t = 1_{\{s_t^{t+\ell-1} = w\}}$ and $\nu(w) = \mathbb{E}[J_0]$. Then
\[
\Var\Bigl(\sum_{t=0}^{N'-1} J_t\Bigr)
\le
N' \,\nu(w) \left(2\ell - 1 + \frac{2\kappa}{1-\kappa}\right).
\]
\end{lemma}
\begin{proof}
$\Var(J_t) \le \nu(w)$. For $1 \le s < \ell$, $|\Cov(J_0,J_s)| \le \nu(w)$. For $s \ge \ell$, the second window is a function of the chain from time $s$ onward, and the total variation between its conditional law given $s_{\ell-1}$ and its stationary law is at most $\kappa^{s-\ell+1}$, because path-space total variation on $\ell$ steps is at most the total variation of the law of the first of those steps. Hence
\[
|\Cov(J_0, J_s)|
=
\bigl|\mathbb{E}\bigl[J_0\bigl(\mathbb{P}(J_s=1 \mid s_{\ell-1}) - \nu(w)\bigr)\bigr]\bigr|
\le
\nu(w)\, \kappa^{s-\ell+1}.
\]
Summing the geometric tail $\sum_{q \ge 1} \kappa^q = \kappa/(1-\kappa)$ and using $N'-s \le N'$ produces the bracket $1 + 2(\ell-1) + 2\kappa/(1-\kappa) = 2\ell-1 + 2\kappa/(1-\kappa)$.
\end{proof}

\begin{theorem}[Expectation and tail]\label{thm:tail}
Let
\[
B_N = \frac{D}{2}\sqrt{\frac{m^\ell}{N'}\left(2\ell - 1 + \frac{2\kappa}{1-\kappa}\right)}.
\]
Then $\mathbb{E}[W_1(\hat\mu_{N,\ell}, \nu_\ell)] \le B_N$. For every $t > 0$,
\[
\mathbb{P}\bigl(W_1 - \mathbb{E}[W_1] \ge t\bigr)
\le
\exp\left(
-\frac{t^2 \, N'^2 (1-\kappa)^2}{2 N D^2 \ell^2}
\right),
\]
and therefore
\[
\mathbb{P}\bigl(\ES \ge (1+\eta)(B_N + t)\bigr)
\le
\exp\left(
-\frac{t^2 \, N'^2 (1-\kappa)^2}{2 N D^2 \ell^2}
\right).
\]
\end{theorem}
\begin{proof}
Unbiasedness $\mathbb{E}[\hat\mu(w)] = \nu(w)$ is stationarity of each window. By Lemma~\ref{lem:var} and Cauchy--Schwarz,
\[
\mathbb{E}\|\hat\mu-\nu\|_1
\le
\sum_w \sqrt{\Var(\hat\mu(w))}
\le
\sqrt{\frac{m^\ell}{N'}\left(2\ell-1+\frac{2\kappa}{1-\kappa}\right)},
\]
where $\sum_w \sqrt{\nu(w)} \le \sqrt{m^\ell}$. Multiply by $D/2$ to pass from $\ell_1$ to $W_1$. That is $B_N$.

For the tail, Lemma~\ref{lem:lip} and Lemma~\ref{lem:influence} give Doob increments bounded by $c/(1-\kappa)$ with $c = D\ell/N'$. Azuma's inequality on $N$ increments is the exponential bound. The comparison $\ES \le (1+\eta) W_1$ lifts it to $\ES$. Replacing the unknown mean by the larger number $B_N$ only enlarges the threshold, so the probability statement remains valid.
\end{proof}

\begin{corollary}[Estimated baseline]\label{cor:estimated}
Let $\hat\nu$ be the $\ell$-marginal of an independent stationary corpus of length $N_0$, with flux Laplacian $\hat L$, algebraic connectivity $\hat\lambda_2$, and unit Fiedler vector $\hat\psi$ chosen with $\langle \hat\psi, \psi \rangle \ge 0$. The triangle inequality gives
\[
W_1(\hat\mu, \nu) \le W_1(\hat\mu, \hat\nu) + W_1(\hat\nu, \nu),
\]
where the second summand obeys Theorem~\ref{thm:tail} with $N_0$ in place of $N$.

For the Fiedler penalty, let $g = \min(\lambda_2,\, \lambda_3 - \lambda_2) > 0$ be the spectral gap of $L$ when $\lambda_2$ is simple (or $\min(\lambda_2,\, \lambda_{2+m_2}-\lambda_2)$ for multiplicity $m_2$ using the basis-independent projection $\Psi\Psi^\top$ of Section~\ref{sec:graph}). Writing $E = \hat L - L$, each undirected edge weight changes by at most the sum of the two directed cylinder errors, so
\[
\|E\|_{\op} \le \|E\|_1 \le 4\,\|\hat\nu - \nu\|_1.
\]
By Weyl's inequality, $|\hat\lambda_2 - \lambda_2| \le \|E\|_{\op}$. Provided $\|E\|_{\op} < g/2$, the Davis--Kahan $\sin\theta$ theorem gives
\[
\|\hat\psi - \psi\|_2 \le \sqrt{2}\,\sin\theta(\hat\psi, \psi) \le \frac{2\sqrt{2}\,\|E\|_{\op}}{g}.
\]
For any transition pair $(u,v)$, the identity $a^2 - b^2 = (a-b)(a+b)$ and Cauchy--Schwarz on the incidence vector $e_u - e_v$ yield
\[
\bigl|(\hat\psi_u - \hat\psi_v)^2 - (\psi_u - \psi_v)^2\bigr|
\le
\|\hat\psi - \psi\|_2 \sqrt{2} \cdot \|\hat\psi + \psi\|_2 \sqrt{2}
\le
4\,\|\hat\psi - \psi\|_2.
\]
Assuming $\supp(\hat\nu) = \supp(\nu)$ (after corpus cover of the allowed transitions), the empirical energies $\hat\Delta_N(\hat\psi)$ and $\hat\Delta_N(\psi)$ on any test window therefore differ by at most $4\|\hat\psi - \psi\|_2$. Since $x \mapsto (1 - x/a)_+$ is $1/a$-Lipschitz in $x$ and $\hat\Delta_N \le 2$, the estimated defect and true defect satisfy
\[
\bigl|\widehat{\defect}(\hat\Delta_N(\hat\psi)) - \defect(\hat\Delta_N(\psi))\bigr|
\le
\frac{4\,\|\hat\psi - \psi\|_2 + \frac{\hat\Delta_N(\psi)}{\lambda_2}|\hat\lambda_2 - \lambda_2|}{\hat\lambda_2}
\le
\frac{\|E\|_{\op}}{\hat\lambda_2}\left(\frac{8\sqrt{2}}{g} + \frac{2}{\lambda_2}\right).
\]
Combined with $\|\hat\nu - \nu\|_1$, which is bounded in expectation by $2 B_{N_0}/D$ and in tail by Lemma~\ref{lem:lip} and Azuma on $N_0$, both the transport shift and the spectral shift vanish at rate $O(N_0^{-1/2})$ once $N_0$ exceeds the cover time of allowed transitions.
\end{corollary}

\begin{remark}[$m$-step contraction]\label{rem:m-step}
If $\kappa(P) = 1$ but $\kappa(P^m) = \kappa_m < 1$, submultiplicativity of $\kappa$ gives $\|\delta_x P^{t} - \delta_y P^{t}\|_{\TV} \le \kappa_m^{\lfloor t/m \rfloor}$. The increment bound becomes $c \, m / (1-\kappa_m)$, and the exponent in Theorem~\ref{thm:tail} replaces $(1-\kappa)$ by $(1-\kappa_m)/m$. H\"older Gibbs states on a mixing subshift of finite type fall under this remark rather than under one-step Doeblin.
\end{remark}

\begin{remark}[Sharpness]\label{rem:sharpness}
The exponent is of order $N t^2 / (D \ell /(1-\kappa))^2$, which is the right shape, but Azuma's constant is far from the variance proxy. Section~\ref{sec:worked} records a chain where $B_N + t$ at a conventional $5$~percent level still exceeds the diameter scale for a few hundred samples, and only drops below half the alternation score near $N = 57000$. Sharper Markov transportation inequalities exist; the argument here does not depend on their constants. The bound is a proof for every finite $N$. It is a usable threshold only once $N$ is large enough that $(1+\eta)(B_N+t)$ sits under the alternative of interest.

A non-stationary start at total-variation distance $\varepsilon_0$ from $q$ shifts each window marginal by at most $\kappa^{t_0}\varepsilon_0$ after $t_0$ burned symbols, which shifts $W_1$ by at most $D \kappa^{t_0}\varepsilon_0$. Burn a logarithmic prefix and the stationary theorem applies to the remainder.
\end{remark}

\begin{remark}[Joint concentration of $W_1$ and $\hat\Delta_N$]\label{rem:joint}
Theorem~\ref{thm:tail} bounds the prefactor by its deterministic worst case $1+\eta$, which multiplies $B_N$ by $1+\eta$ even though $\mathbb{E}_\nu[\hat\Delta_N] = \lambda_2$. Because each legal transition jump satisfies $(\psi_u - \psi_v)^2 \le 2\|\psi\|_2^2 = 2$, changing one symbol $s_j$ alters at most $\ell$ terms of $\hat\Delta_N$, shifting $\hat\Delta_N$ by at most $c_\Delta = 2\ell/N'$. By Lemma~\ref{lem:influence} and Azuma's inequality on $\lambda_2 - \hat\Delta_N$, for every $s \in (0,\lambda_2]$,
\[
\mathbb{P}\bigl(\defect(\hat\Delta_N) \ge s/\lambda_2\bigr)
=
\mathbb{P}\bigl(\lambda_2 - \hat\Delta_N \ge s\bigr)
\le
\exp\left(
-\frac{s^2 \, N'^2 (1-\kappa)^2}{8 N \ell^2}
\right).
\]
By the union bound, with probability at least $1 - \delta_1 - \delta_2$,
\[
\ES
\le
(B_N + t)\left(1 + \frac{\eta s}{\lambda_2}\right),
\]
where $t$ and $s$ are the Azuma deviations at levels $\delta_1$ and $\delta_2$. Since $B_N$ and $s$ both decay as $O(N^{-1/2})$, the leading threshold term is $B_N + O(N^{-1/2})$ rather than $(1+\eta)B_N$, reducing the sample size required to separate a fixed alternative by a factor of up to $(1+\eta)^2$.
\end{remark}

\section{Worked Chain}\label{sec:worked}

\subsection{Baseline}

\[
P = \begin{pmatrix} 1/10 & 4/5 & 1/10 \\ 2/5 & 1/5 & 2/5 \\ 1/10 & 4/5 & 1/10 \end{pmatrix},
\qquad
q = \bigl(\tfrac14,\tfrac12,\tfrac14\bigr).
\]
The row $q$ is stationary by direct multiplication. Symbol entropy is
\[
-\tfrac12 \log_2 \tfrac14 - \tfrac12 \log_2 \tfrac12 = \tfrac32 \ \text{bits}.
\]
The entropy rate $-\sum q_i P_{ij} \log_2 P_{ij}$ evaluates to $1.2219280948873625$~bits in binary64. One-step Dobrushin coefficient: rows $0$ and $2$ agree, and
\[
\|P(0,\cdot)-P(1,\cdot)\|_{\TV} = \tfrac12(0.3+0.6+0.3) = \tfrac35.
\]
So $\kappa = 3/5$. Ground metric on symbols $\delta(i,j) = |i-j|$, on pairs $d = |i-k| + \tfrac12 |j-l|$, so $D = 3$. Penalty weight $\eta = 1$.

Fluxes of the off-diagonal edges: $w(0,1) = w(1,2) = 2/5$, $w(0,2) = 1/20$. The weighted Laplacian is
\[
L = \begin{pmatrix}
9/20 & -2/5 & -1/20 \\
-2/5 & 4/5 & -2/5 \\
-1/20 & -2/5 & 9/20
\end{pmatrix}.
\]
The unit vector $\psi = (-1/\sqrt{2},\, 0,\, 1/\sqrt{2})$ satisfies $L\psi = \tfrac12 \psi$ by multiplication, and $\psi \perp \mathbf{1}$. The quadratic form reproduces the same eigenvalue:
\[
\frac25 \cdot \frac12 + \frac25 \cdot \frac12 + \frac{1}{20} \cdot 2 = \frac12.
\]
So $\lambda_2 = \Delta_* = 1/2$. The trace of $L$ is $17/10$, so the remaining eigenvalue is $17/10 - 1/2 = 6/5$. The eigenvalue $\lambda_2$ is simple. Squared jumps: $(\psi_0-\psi_1)^2 = 1/2$, $(\psi_0-\psi_2)^2 = 2$, self-loops $0$.

\subsection{Five population alternatives}

The edge-Lipschitz criterion is used below. The metric $d$ is the shortest-path metric of the product-of-paths graph whose horizontal edges (first symbol changes by $1$) have length $1$ and whose vertical edges (second symbol changes by $1$) have length $1/2$. A function on that graph is $1$-Lipschitz if and only if each edge changes the function by at most the edge length; the inequality along a general pair follows by summing a shortest path.

\medskip\noindent\textbf{Constant state $1$.} The empirical pair measure is $\delta_{(1,1)}$. Because the source is a Dirac, there is only one coupling, and
\begin{align*}
W_1
&= \sum_{i,j} \nu(i,j)\, d\bigl((1,1),(i,j)\bigr) \\
&= 4\cdot\frac{1}{40}\cdot\frac32 + 2\cdot\frac15\cdot 1 + 2\cdot\frac15\cdot\frac12 \\
&= \frac{3}{20} + \frac25 + \frac15
= \frac34.
\end{align*}
The four corner words $(0,0),(0,2),(2,0),(2,2)$ each have mass $1/40$ and distance $3/2$; $(0,1)$ and $(2,1)$ each have mass $1/5$ and distance $1$; $(1,0)$ and $(1,2)$ each have mass $1/5$ and distance $1/2$. Energy $\hat\Delta = 0$, defect $= 1$, prefactor $= 2$, so $\ES = 3/2$. Entropy rate $0$.

\medskip\noindent\textbf{Alternation $0,2,0,2,\ldots$.} Occupation $\mu(0,2) = \mu(2,0) = 1/2$. Energy $= 2 > \lambda_2$, so the one-sided defect is $0$ and $\ES = W_1$.

Dual witness, rows indexed by the first symbol and columns by the second:
\[
f = \begin{pmatrix}
0 & 1/2 & 1 \\
0 & -1/2 & 0 \\
1 & 1/2 & 0
\end{pmatrix}.
\]
Adjacent horizontal differences are in $\{0,1\}$ and adjacent vertical differences are all $1/2$, so $f$ is $1$-Lipschitz. With $\nu(0,0)=\nu(0,2)=\nu(2,0)=\nu(2,2)=1/40$, $\nu(0,1)=\nu(1,0)=\nu(1,2)=\nu(2,1)=1/5$, $\nu(1,1)=1/10$,
\begin{align*}
\sum f\,d(\mu-\nu)
&= 2\cdot\bigl(\tfrac12 - \tfrac{1}{40}\bigr)\cdot 1
+ 2\cdot\bigl(-\tfrac15\bigr)\cdot\tfrac12
+ \bigl(-\tfrac{1}{10}\bigr)\cdot\bigl(-\tfrac12\bigr) \\
&= \tfrac{19}{20} - \tfrac15 + \tfrac{1}{20}
= \tfrac45.
\end{align*}
So $W_1 \ge 4/5$. The coupling that moves the $\mu$-mass at $(0,2)$ onto $\nu$ by the amounts
\[
(0,0):\tfrac{1}{40},\quad
(0,1):\tfrac15,\quad
(0,2):\tfrac{1}{40},\quad
(1,1):\tfrac{1}{20},\quad
(1,2):\tfrac15,
\]
and the symmetric amounts from $(2,0)$ onto $(1,0),(1,1),(2,0),(2,1),(2,2)$, has both correct marginals (each $\nu$-atom is hit once, and each source has mass $1/2$). Its cost is
\[
2\left(\frac{1}{40} + \frac{1}{10} + \frac{3}{40} + \frac15\right) = \frac45.
\]
Hence $W_1 = 4/5$ and $\ES = 4/5$. Entropy rate $0$.

The two-sided ratio of Section~\ref{subsec:two-sided} multiplies this by $\lambda_2(1+\varepsilon)/(\hat\Delta+\varepsilon\lambda_2)$ at $\varepsilon=1/4$:
\[
\frac{(1/2)\cdot(5/4)}{2 + (1/4)\cdot(1/2)} \cdot \frac45 = \frac{4}{17}.
\]

\medskip\noindent\textbf{Fair noise on the cut $\{0,2\}$.} i.i.d.\ uniform on $\{0,2\}$ has entropy rate $1$~bit. Pair masses $1/4$ on $(0,0),(0,2),(2,0),(2,2)$. Energy $= 2\cdot(1/4)\cdot 2 = 1 > 1/2$, defect $0$.

Dual witness
\[
f = \begin{pmatrix}
0 & -1/2 & 0 \\
-1 & -3/2 & -1 \\
0 & -1/2 & 0
\end{pmatrix}
\]
has adjacent differences inside the edge lengths, so it is $1$-Lipschitz. Then $\mu(f) = 0$ and
\[
\nu(f)
= 2\cdot\frac15\cdot\bigl(-\tfrac12\bigr)
+ 2\cdot\frac15\cdot(-1)
+ \frac{1}{10}\cdot\bigl(-\tfrac32\bigr)
= -\frac{3}{4},
\]
so $(\mu-\nu)(f) = 3/4$. The matching coupling, written as source, target, mass, is
\begin{align*}
&(0,0)\to(0,0),\ \tfrac{1}{40};
&&(0,0)\to(0,1),\ \tfrac{7}{40};
&&(0,0)\to(1,1),\ \tfrac{1}{20};\\
&(0,2)\to(0,1),\ \tfrac{1}{40};
&&(0,2)\to(0,2),\ \tfrac{1}{40};
&&(0,2)\to(1,2),\ \tfrac15;\\
&(2,0)\to(1,0),\ \tfrac15;
&&(2,0)\to(2,0),\ \tfrac{1}{40};
&&(2,0)\to(2,1),\ \tfrac{1}{40};\\
&(2,2)\to(1,1),\ \tfrac{1}{20};
&&(2,2)\to(2,1),\ \tfrac{7}{40};
&&(2,2)\to(2,2),\ \tfrac{1}{40}.
\end{align*}
Row masses are $1/4$. Column masses reproduce $\nu$. The cost sums to
\[
\frac{7}{80}+\frac{3}{40}+\frac{1}{80}+\frac15+\frac15+\frac{1}{80}+\frac{3}{40}+\frac{7}{80} = \frac34.
\]
Hence $W_1 = 3/4$ and $\ES = 3/4$.

\medskip\noindent\textbf{Fair noise on one side of the cut, $\{0,1\}$.} Entropy rate again $1$~bit. Energy:
\[
\frac14 \cdot \frac12 + \frac14 \cdot \frac12 = \frac14,
\]
so the defect is $(1/2 - 1/4)/(1/2) = 1/2$ and the prefactor is $3/2$.

The function $f(i,j) = -(i + j/2)$ obeys
\[
|f(i,j)-f(k,l)| = \bigl|(i-k) + \tfrac12(j-l)\bigr| \le |i-k| + \tfrac12|j-l|.
\]
It is $1$-Lipschitz. Then
\[
\mu(f) = \tfrac14\bigl(0 - \tfrac12 - 1 - \tfrac32\bigr) = -\tfrac34,
\]
and, summing $\nu(i,j)f(i,j)$ over $(0,1),(0,2),(1,0),(1,1),(1,2),(2,0),(2,1),(2,2)$,
\[
\nu(f)
= -\frac{1}{10}-\frac{1}{40}-\frac15-\frac{3}{20}-\frac25-\frac{1}{20}-\frac12-\frac{3}{40}
= -\frac{60}{40}
= -\frac32,
\]
so $(\mu-\nu)(f) = 3/4$. A feasible coupling of cost $3/4$ is
\begin{align*}
&(0,0)\to(0,0),\ \tfrac{1}{40};
&&(0,0)\to(0,2),\ \tfrac{1}{40};
&&(0,0)\to(1,2),\ \tfrac15;\\
&(0,1)\to(0,1),\ \tfrac15;
&&(0,1)\to(2,1),\ \tfrac{1}{20};\\
&(1,0)\to(1,0),\ \tfrac15;
&&(1,0)\to(2,0),\ \tfrac{1}{40};
&&(1,0)\to(2,2),\ \tfrac{1}{40};\\
&(1,1)\to(1,1),\ \tfrac{1}{10};
&&(1,1)\to(2,1),\ \tfrac{3}{20}.
\end{align*}
Row masses are $1/4$ and column masses reproduce $\nu$. The cost is
\[
\frac{1}{40}+\frac25+\frac{1}{10}+\frac{1}{40}+\frac{1}{20}+\frac{3}{20} = \frac34.
\]
Hence $W_1 = 3/4$ and $\ES = (3/4)\cdot(3/2) = 9/8$.

\medskip\noindent\textbf{Local jitter at scale $1/10$.} Each symbol is retained with probability $9/10$ and moved to a neighbor otherwise, the two neighbors of state $1$ splitting the mass $1/10$. The transition kernel $K$ satisfies $qK = q$ and $K\psi = \tfrac{9}{10}\psi$. Because the pushed $1$-marginal is still $q$, the pushed symbol entropy is still $3/2$~bits and $\mathbb{E}[\psi(S'_t)^2] = \tfrac14$. Under $\nu$, the four corner pairs $(0,0),(2,2),(0,2),(2,0)$ all have mass $1/40$, so $\mathbb{E}_\nu[\psi(S_0)\psi(S_1)] = 0$ and hence
\[
\Delta(\nu^K)
=
2\,\mathbb{E}[\psi(S'_0)^2] - 2\bigl(\tfrac{9}{10}\bigr)^2 \mathbb{E}_\nu[\psi(S_0)\psi(S_1)]
=
\frac12.
\]
Thus $\Delta(\nu^K) = \lambda_2 = 1/2$ and $\defect(\Delta(\nu^K)) = 0$ hold identically in exact arithmetic. Solving the transport LP on the nine words in exact rational arithmetic returns $W_1(\nu^K_2, \nu_2) = 27/500 = 0.054$ with zero primal--dual gap, so $\ES = 27/500 = 0.054$. It sits an order of magnitude below every structured score in the table.

\subsection{What the table separates}

\begin{table}[ht]
\centering
\caption{Population scores on the $3$-state worked chain ($\ell=2, \eta=1$).}\label{tab:pop}
\begin{tabular}{@{}lcccc@{}}
\toprule
Law & Entropy rate (bits) & $W_1$ & Defect & $\ES$ \\
\midrule
Baseline $\nu$ & $1.222$ & $0$ & $0$ & $0$ \\
Jitter, scale $1/10$ & symbol entropy $3/2$ & $27/500 = 0.054$ & $0$ & $27/500 = 0.054$ \\
i.i.d.\ on the cut $\{0,2\}$ & $1$ & $3/4$ & $0$ & $3/4$ \\
Alternation $0 \leftrightarrow 2$ & $0$ & $4/5$ & $0$ & $4/5$ \\
i.i.d.\ on one side $\{0,1\}$ & $1$ & $3/4$ & $1/2$ & $9/8$ \\
Constant state $1$ & $0$ & $3/4$ & $1$ & $3/2$ \\
\bottomrule
\end{tabular}
\end{table}

The cut-confined process and the one-sided process have the same entropy rate and the same transport cost. The one that fails to pay the Fiedler energy scores $9/8$ rather than $3/4$. That is the whole content of the penalty: among equal-cost deviations, loss of algebraic connectivity ranks higher. The constant state is the extreme case, prefactor $2$.

The alternation pays \emph{more} than the baseline energy, because the rare edge is the cut itself. The one-sided defect stays $0$ and the score remains the transport cost $4/5$. A ratio that treats extra energy as a credit shrinks this orbit to $4/17$.

\subsection{Finite windows}

Eight hundred independent stationary windows of length $N = 400$, generator seed $7$, scored with the one-sided definition:

\begin{table}[ht]
\centering
\caption{Finite-window Monte Carlo scores ($800$ windows, $N = 400$, seed $7$).}\label{tab:mc}
\begin{tabular}{@{}lcccc@{}}
\toprule
Source & Mean & 95th & 99th & Max \\
\midrule
Baseline & $0.0511$ & $0.1069$ & $0.1479$ & $0.2356$ \\
Jitter, scale $1/10$ & $0.0769$ & $0.1143$ & $0.1565$ & $0.2187$ \\
i.i.d.\ on the cut & $0.750$ & $0.750$ & $0.750$ & $0.750$ \\
Alternation & $0.800$ & $0.800$ & $0.800$ & $0.800$ \\
i.i.d.\ on one side & $1.123$ & $1.220$ & $1.264$ & $1.339$ \\
Constant state & $1.500$ & $1.500$ & $1.500$ & $1.500$ \\
\bottomrule
\end{tabular}
\end{table}

Every structured window in the sample sits above the worst baseline window. Jitter raises the mean from $0.051$ to $0.077$ and does not enter the structured range. This is a Monte Carlo illustration. It is not a substitute for Theorem~\ref{thm:tail}.

A second sample of $800$ baseline windows at the same $N$ and seed, scored with the two-sided ratio at $\varepsilon = 1/4$, had mean $0.0502$, 99th percentile $0.145$, and maximum $0.202$. The alternation's two-sided score $4/17 \approx 0.235$ is above that maximum by about $0.033$. The one-sided gap, against a baseline maximum of $0.236$, is $0.564$.

\subsection{The explicit threshold on this chain}

Here $m^\ell = 9$, $\ell = 2$, $D = 3$, $\kappa = 3/5$, bracket $= 3 + 2\cdot(3/5)/(2/5) = 6$, $\eta = 1$. Theorem~\ref{thm:tail} at $N = 400$ ($N' = 399$) produces an expectation proxy near $0.55$ and a one-sided $5$~percent Azuma deviation near $1.85$. The resulting threshold $(1+\eta)(B_N+t)$ is near $4.8$, which is larger than every score in the table ($W_1 \le 3$, so $\ES \le 6$). The inequality is true and too loose to test at this $N$.

The same closed form, evaluated on a grid of sample sizes, first falls below $2/5$ (half the alternation score) at $N = 57000$. Using the joint union bound of Remark~\ref{rem:joint} (splitting the $0.05$ level as $\delta_1 = 0.04$ on $W_1$ and $\delta_2 = 0.01$ on $\hat\Delta_N$) removes the factor of $2$ on $B_N$ and drops the same $2/5$ crossing from $N = 57000$ down to $N \approx 16400$. From that $N$ on, the explicit test ``flag when $\ES$ exceeds the concentration threshold at level $0.05$'' has false-alarm probability at most $0.05$ under the baseline and fires on the alternation, whose score is $4/5$ for every $N$. The Monte Carlo separation at $N = 400$ is real in the sample and is not yet the separation proved by the Azuma constant.

The empirical mean of $W_1$ on baseline windows was about $0.050$, against an expectation proxy $B_N \approx 0.55$. The gap is the price of Cauchy--Schwarz over all nine words and of Azuma's quadratic tail. It is recorded so the proxy is not mistaken for a typical value.

\section{Reading for Event Telemetry}\label{sec:telemetry}

The observation is a sequence of bins. For hardware performance counters, a time bucket produces a vector of counter deltas and the bin is a cell of a fixed quantizer fit on the trusted corpus. For packet or flow telemetry, the bin is a cell of inter-arrival time, size, or direction, again fixed from the corpus. The symbol clock has to be chosen on purpose.
\begin{itemize}
\item If one symbol is emitted per event, a pure speed change that preserves the sequence of event types does not change the symbolic window, and $\ES$ is invariant under that reclocking. Volume, in the sense of events per second, is invisible.
\item If one symbol is emitted per fixed clock and the bin includes a count, volume enters the symbol and can move $W_p$.
\end{itemize}

The baseline corpus estimates $q$, $P$, the flux Laplacian, and one Fiedler vector. A sliding window then maintains $\ell$-word counts, solves or approximates $W_p$ against the corpus marginal, and evaluates $\hat\Delta_N$ as an average of $(\psi_u-\psi_v)^2$ on legal transitions. Algebraic connectivity itself is computed once.

What a large value means: the window's transition tensor is an expensive transport away from the corpus marginal, or it is a moderate transport that fails to carry the corpus Fiedler energy. In the second case the occupied transitions sit inside a region the Fiedler coordinate does not cross, which is the spectral form of a bottleneck in the allowed grammar. Ordinary high-entropy motion inside the corpus geometry stays cheap by Proposition~\ref{prop:jitter}, up to the factor $1+\eta$.

A low-and-slow deviation whose binned transitions put mass $\alpha$ at symbol-distance $d$ from the corpus bulk has $W_1$ of order $\alpha d$, hence $\ES$ of that order at least, even when the scalar entropy rate of the window remains within a few tenths of a bit of the corpus rate. The cut-confined row of the table is a miniature of that pattern: entropy rate $1$ against a baseline rate $1.222$, score $3/4$ against a baseline finite-window 99th percentile near $0.15$.

The computable regime is a small alphabet. Exact $W_1$ is a linear program with $(m^\ell)^2$ variables. For $m$ in the tens and $\ell = 2$ or $3$, an entropic approximation (Sinkhorn) is the practical evaluation of the same definition. The Fiedler vector is an $m^k \times m^k$ eigenproblem, once per corpus.

\section{Limits}\label{sec:limits}

The zero set is exactly the set of windows whose $\ell$-marginal equals $\nu_\ell$. A modification that preserves $\ell$-word frequencies, and therefore also preserves $\hat\Delta_N$, has population score zero. Matching the physical marginal is a blind spot of this statistic. That is the content of Proposition~\ref{prop:zero}, not a procedure.

A coarser partition quotients the symbol metric, so distances and $W_p$ fall. A sub-manifold that lives inside a single bin is invisible. The partition is part of the definition.

$\nu$ estimated from a short or non-representative corpus moves the target. Corollary~\ref{cor:estimated} is the corresponding error term. If the trusted corpus already mixes several regimes, the reference $\nu$ is a mixture, and a window that sits in one regime can score as a deviation from the mixture.

The Azuma threshold is explicit and conservative. On the chain of Section~\ref{sec:worked} it becomes a separating test near $N = 57000$ (or $N \approx 16400$ under the joint bound of Remark~\ref{rem:joint}) at the $5$~percent level, while samples already separate at a few hundred symbols. Using the Azuma proxy as an operational alarm at moderate $N$ will not fire on the alternatives in the table, because the threshold lies above them.

No SRB theorem is claimed for kernels, protocol state machines, or counter traces. For those, $\ES$ is a comparison of a finite symbolic window with a chosen reference measure, plus a shortfall against the algebraic connectivity of that reference's allowed graph.

\section{Computation Checklist}\label{sec:checklist}

\begin{enumerate}
\item Fix the quantizer and the block length $\ell$. Fit it on the corpus only.
\item Estimate the $\ell$-marginal $\hat\nu$ and the flux weights $w$. Build $L$. Compute $\lambda_2$ and $\psi$.
\item On each window, form $\hat\mu$ and $\hat\Delta_N$ (forbidden transitions contribute $0$ to the energy).
\item Compute $W_p(\hat\mu, \hat\nu)$ and multiply by $1 + \eta (\lambda_2 - \hat\Delta_N)_+/\lambda_2$.
\item Compare with the Theorem~\ref{thm:tail} / Remark~\ref{rem:joint} threshold if $N$ is large enough that the threshold lies under the deviation of interest, and add the corpus term from Corollary~\ref{cor:estimated}. At moderate $N$, treat the threshold as a proved sufficient bound, and calibrate an operational cutoff on held-out baseline windows. Report that calibration as empirical.
\end{enumerate}

\begin{thebibliography}{99}

\bibitem{Azuma1967}
K.~Azuma, \emph{Weighted sums of certain dependent random variables}, Tohoku Math.\ J.\ \textbf{19} (1967), 357--367.

\bibitem{Bowen1975}
R.~Bowen, \emph{Equilibrium States and the Ergodic Theory of Anosov Diffeomorphisms}, Lecture Notes in Mathematics, vol.~470, Springer, 1975.

\bibitem{DavisKahan1970}
C.~Davis and W.~M.~Kahan, \emph{The rotation of eigenvectors by a perturbation.~III}, SIAM J.\ Numer.\ Anal.\ \textbf{7} (1970), 1--46.

\bibitem{Fiedler1973}
M.~Fiedler, \emph{Algebraic connectivity of graphs}, Czechoslovak Math.\ J.\ \textbf{23} (1973), 298--305.

\bibitem{GrassbergerProcaccia1983}
P.~Grassberger and I.~Procaccia, \emph{Characterization of strange attractors}, Phys.\ Rev.\ Lett.\ \textbf{50} (1983), 346--349.

\bibitem{MuskulusVerduynLunel2011}
M.~Muskulus and S.~Verduyn-Lunel, \emph{Wasserstein distances in the analysis of time series and dynamical systems}, Physica~D \textbf{240} (2011), 45--58.

\bibitem{Pesin1977}
Ya.~B.~Pesin, \emph{Characteristic Lyapunov exponents and smooth ergodic theory}, Russian Math.\ Surveys \textbf{32} (1977), no.~4, 55--114.

\bibitem{Riekert2022}
A.~Riekert, \emph{Convergence rates for empirical measures of Markov chains in dual and Wasserstein distances}, Statist.\ Probab.\ Lett.\ \textbf{189} (2022), 109577.

\bibitem{Ruelle1976}
D.~Ruelle, \emph{A measure associated with Axiom~A attractors}, Amer.\ J.\ Math.\ \textbf{98} (1976), 619--654.

\bibitem{Sinai1972}
Ya.~G.~Sinai, \emph{Gibbs measures in ergodic theory}, Russian Math.\ Surveys \textbf{27} (1972), no.~4, 21--69.

\bibitem{ZivLempel1978}
J.~Ziv and A.~Lempel, \emph{Compression of individual sequences via variable-rate coding}, IEEE Trans.\ Inform.\ Theory \textbf{24} (1978), 530--536.

\end{thebibliography}

\end{document}